\documentclass[12pt]{article}
\usepackage[utf8]{inputenc}
\usepackage[T1]{fontenc}
\usepackage{hyperref}
\title{Scots Language Resources for K-12 Curriculum Development}
\author{Compiled by Claude Sonnet 5.5 (claude-sonnet-5-5) for Andr\'as Kornai}
\date{}
\begin{document}
\maketitle

\section{Introduction}
Scots (ISO 639-3 sco) is a Germanic language of Scotland and Ulster. It is well documented and has some school presence, mostly as a cross-curricular topic or as an ``L3'' under Scotland's 1+2 languages approach. This dossier lists only resources that appeared in web searches.

\section{Primary Educational Resources}
\begin{itemize}
\item Education Scotland Scots Language hub: an animated history of Scots, and materials for primary and senior phase. It also covers the SQA Scots Language Award. \url{https://www.scotslanguage.com/articles/view/id/4641/type}
\item ``New online resource for teaching the Scots language'' (SCILT). \url{https://scilt.org.uk/News/NewsView/tabid/1311/articleType/ArticleView/articleId/6168/New-online-resource-for-teaching-the-Scots-language.aspx}
\item ``Teaching Scots as part of 1+2 languages'' (Scots Language Centre). \url{https://www.scotslanguage.com/articles/view/id/6241}
\item Wee Windaes Teaching Resources (National Library of Scotland, with Scots Scriever Hamish MacDonald), for primary and secondary literacy. \url{https://www.scotslanguage.com/articles/node/id/847/type}
\item NLS announcement of Wee Windaes, launched 1 December 2016. \url{https://media.nls.uk/news/website-offers-windae-into-scots-language}
\item Education news from the Scots Language Centre (Winter 2017). \url{https://www.scotslanguage.com/articles/node/id/603}
\end{itemize}

\section{Linguistic and Reference Materials}
\begin{itemize}
\item Dictionary of the Scots Language, combining the Scottish National Dictionary and the Dictionary of the Older Scottish Tongue; free online. \url{https://scotlex.org/dictionary-of-the-scots-language/}
\item Dictionaries of the Scots Language, overview. \url{https://en.wikipedia.org/wiki/Dictionaries_of_the_Scots_Language}
\item Dictionary of the Scots Language, overview. \url{https://en.wikipedia.org/wiki/Dictionary_of_the_Scots_Language}
\item DSL publications. \url{https://dsl.ac.uk/?p=14}
\item DSL preface (PDF). \url{https://scotlex.org/wp-content/uploads/2020/06/dslpreface.pdf}
\end{itemize}

\section{Religious and Community Resources}
\begin{itemize}
\item W. L. Lorimer, \emph{The New Testament in Scots}; the Gospel of John is published by Canongate. \url{https://canongate.co.uk/books/1931-johns-gospel-from-the-new-testament-in-scots-translated-by-william-laughton-lorimer/}
\item Luke's Gospel in Scots (Canongate). \url{https://canongate.co.uk/books/1930-lukes-gospel-from-the-new-testament-in-scots-translated-by-william-laughton-lorimer/}
\item Scots Language Centre (scotslanguage.com), the main community and advocacy site. \url{https://www.scotslanguage.com/articles/view/id/4308}
\end{itemize}

\section{Digital and Audio Resources}
\begin{itemize}
\item Audiobook readings of Lorimer's Gospels by Tom Fleming. \url{https://audible.com/topic/audiobooks-scots-bible}
\item Wee Windaes site texts with audio and video (digitised books at NLS). \url{https://digital.nls.uk/selected-wee-windaes-books/archive/127591362?mode=zoom}
\item SQA/Qualifications Scotland blog on Scots language. \url{https://blogs.sqa.org.uk/scottishstudies/tag/scots-language}
\end{itemize}

\section{Curriculum Development Context}
The materials found cover cultural and literary content, an SQA award and teaching resources for primary and secondary pupils. No complete Scots-medium K-12 curriculum or graded course textbooks were found. The searches did not locate current Scottish Government policy text, so that should be checked. Scots has no fixed standard spelling, so a curriculum would need to choose a house orthography, and the DSL and Wee Windaes texts offer a base for that.
\end{document}